\reading{LTR-10}{Liquidity Risk Reporting and Stress Testing}{STRIKE FIRST}{2}

\begin{lrkeybox}[Learning objectives — official 2026 spine wording]
\begin{enumerate}[label=\textbf{\alph*.}]
\item Describe best practices for the reporting of a bank's liquidity position.
\item Compare and interpret different types of liquidity risk reports.
\item Explain the process of reporting a liquidity stress test and interpret a liquidity stress test report.
\end{enumerate}
{\footnotesize\color{FRMGrey}Source: Moorad Choudhry, \textit{The Principles of Banking}, Chapter 14.}
\end{lrkeybox}

\begin{lrtrapbox}[This is the thinnest reading in the source notes]
LTR-10 exists in the Crash layer only, and objectives \textbf{a} and \textbf{b} together occupy four sentences --- one of which trails off mid-colon and one of which names a report and then stops. Both are gap-filled below. This is a \textbf{STRIKE FIRST} reading, so the thinness sat exactly where it could do the most damage.
\end{lrtrapbox}

% =====================================================================
\lo{LTR-10 a}{Describe best practices for the reporting of a bank's liquidity position.}

A bank will produce a number of liquidity reports in the normal course of business, on a \textbf{daily, weekly, monthly and quarterly} basis. It is important that the format of liquidity management information is both \textbf{transparent and accessible}.

\subsection{What the source notes carry}

\begin{itemize}
\item \textbf{FSA requirements.} The FSA sets liquidity reporting frequency for banks in the UK. \textbf{Daily reports} --- cash flow and mismatch analysis --- are crucial for understanding a bank's liquidity.
\item \textbf{Off-balance-sheet items.} Banks may include some off-balance-sheet items in their reporting.
\end{itemize}

\gapbadge
\gapnote{The source notes stop there, mid-sentence on the off-balance-sheet point, and name no report and no frequency. The rest of this objective is written from the GARP curriculum (Choudhry, Chapter 14).}

\begin{lrgapbox}
The objective asks for \emph{best practices}. That means the governing principle (who sets frequency, who reviews), the specific UK reporting schedule the reading uses as its worked illustration, and why off-balance-sheet commitments are reported \emph{separately} rather than folded in.
\end{lrgapbox}

\begin{lrkeybox}[Who stipulates the frequency]
In general, the main liquidity reports are \textbf{required by the regulator, who stipulates their frequency}. The \textbf{ALCO} should view this reporting. Where reporting requirements are waived or modified --- smaller institutions and foreign branches that are not ILAS firms --- the regulatory authority agrees the format and frequency on a \textbf{case-by-case basis}.
\end{lrkeybox}

\subsection{UK liquidity reporting requirements: standard ILAS firms}

\begin{center}
\scriptsize
\begin{longtable}{@{}p{3.2cm} p{5.1cm} p{2.6cm} p{3.3cm}@{}}
\toprule
\rowcolor{FRMSoft}\textbf{Report} & \textbf{Description} & \textbf{Frequency} & \textbf{Submission deadline} \\
\midrule
\endfirsthead
\toprule
\rowcolor{FRMSoft}\textbf{Report} & \textbf{Description} & \textbf{Frequency} & \textbf{Submission deadline} \\
\midrule
\endhead
\textbf{FSA047} \newline Daily Flows &
Daily cash flows out to 3 months; analyses survival period. &
BAU: \textbf{weekly}. \newline Firm-specific and/or market-wide liquidity \textbf{stress: daily}. &
BAU: end-of-day Monday. Stress: end of day following business day. \\
\addlinespace
\textbf{FSA048} \newline Enhanced mismatch report &
ILAS risk drivers and contractual cash flows across the full maturity spectrum. &
As FSA047. & As FSA047. \\
\addlinespace
\textbf{FSA050} \newline Liquidity Buffer qualifying securities &
Granular analysis of the firm's marketable asset holdings. &
Monthly & 15 business days after month-end. \\
\addlinespace
\textbf{FSA051} \newline Funding concentration &
The firm's borrowings of unsecured wholesale funds (excludes primary issuance), by counterparty class. &
Monthly & 15 business days after month-end. \\
\addlinespace
\textbf{FSA052} \newline Wholesale liabilities &
Daily transaction prices and transacted volumes for wholesale unsecured liabilities. &
Weekly & End-of-day Tuesday. \\
\addlinespace
\textbf{FSA053} \newline Retail, SME and large enterprises corporate funding &
The firm's retail and corporate funding profile, and the \textbf{stickiness of retail deposits}. &
Quarterly & 15 business days after month-end. \\
\addlinespace
\textbf{FSA054} \newline Currency analysis &
Analysis of foreign exchange exposures on the firm's balance sheet. &
Quarterly & 15 business days after month-end. \\
\addlinespace
\textbf{Off-balance-sheet report} &
Aggregate \textbf{undrawn committed facilities}. &
Monthly & 15 business days after month-end. \\
\bottomrule
\end{longtable}
\end{center}
\figcap{The frequency column is the examinable structure, not the report codes. Two reports are daily-or-weekly and switch to \emph{daily} under stress; two are weekly or monthly; three are quarterly. The only one that changes frequency when conditions change is the daily-flows pair --- everything else runs on a fixed calendar.}

\begin{lrtrapbox}[FSA047 is the one that accelerates under stress]
Under business as usual FSA047 and FSA048 are \textbf{weekly}; under firm-specific or market-wide stress they become \textbf{daily}. An option describing them as daily at all times, or as unchanged under stress, has corrupted the one conditional in the table.
\end{lrtrapbox}

\begin{lrkeybox}[Why undrawn commitments are reported separately]
In a stress situation a bank can expect \textbf{un-utilised liquidity and funding lines to be drawn down}, as customers experience funding difficulties of their own. The existence of undrawn commitments can therefore \textbf{exacerbate funding shortages at exactly the wrong time}, which is why liquidity metrics include them --- and why they are reported separately rather than netted into the balance sheet view.
\end{lrkeybox}

\begin{lrkeybox}[Qualitative reporting alongside the numbers]
Group treasury also collects \textbf{qualitative} monthly liquidity highlights from branches and subsidiaries, explaining: significant changes in the 1-week and 1-month liquidity ratios; changes to the cash and liquidity gap in the cumulative liquidity model; significant changes to the Liquidity Risk Factor; growth or shrinkage of asset books; changes to the inter-group borrowing and lending position, with counterparties named for large deals; increases or decreases in corporate deposits, with an estimated roll-over confidence level for large dated transactions; changes in retail deposits; and the average daily opening cash position.
\end{lrkeybox}

% =====================================================================
\lo{LTR-10 b}{Compare and interpret different types of liquidity risk reports.}

\gapbadge
\gapnote{The source notes name the deposit tracker report and stop. The remaining report types below are written from the GARP curriculum (Choudhry, Chapter 14).}

\begin{lrgapbox}
The objective asks the candidate to \emph{compare and interpret}, so what matters is what each report is for, how often it runs, and what an adverse reading of it looks like.
\end{lrgapbox}

\begin{center}
\scriptsize
\begin{longtable}{@{}p{3.3cm} p{5.9cm} p{5.0cm}@{}}
\toprule
\rowcolor{FRMSoft}\textbf{Report} & \textbf{What it shows} & \textbf{How to read it} \\
\midrule
\endfirsthead
\toprule
\rowcolor{FRMSoft}\textbf{Report} & \textbf{What it shows} & \textbf{How to read it} \\
\midrule
\endhead
\textbf{Deposit tracker report} &
The current size of deposits together with a \textbf{forecast} of expected future levels. Month-end actuals by customer type, the change from each month-end, aggregate customer assets and hence the \textbf{LTD ratio}, and a forecast to year-end. &
Tracked \textbf{weekly and monthly}, because it gives an idea of the LTD ratio in the immediate short term. The LTD is a key management liquidity ratio. The forecast element rests on \textbf{objective judgement}, informed by the historical trend and inputs from relationship managers. \\
\addlinespace
\textbf{Daily liquidity report} &
A spreadsheet detailing the bank's \textbf{liquid and marketable assets} together with liabilities, up to 1-year maturity and beyond. &
Gives an end-of-day picture of the liquidity position for Treasury and Finance. Each branch and subsidiary completes one; a bank with only a branch structure may \textbf{aggregate} the report. \\
\addlinespace
\textbf{Funding maturity gap (``mismatch'') report} &
The maturity gap per \textbf{time bucket}, for all assets and liabilities, with an adjustment for liquid securities. Includes the cumulative liquidity cash flow. &
The two reports can be combined. The same report generates the \textbf{cash flow survival horizon} report. \\
\addlinespace
\textbf{Funding concentration report} &
Reliance on each source or sector of funds, \textbf{including intra-group funds}. &
Key management information for senior Treasury and relationship managers. The governing principle is \textbf{funding diversity}: a bank should not become over-reliant on a single source. \\
\addlinespace
\textbf{Large depositor concentration report} &
Depositors above a size threshold --- for example USD 50 million, though a bank may define ``large'' as a percentage of total liabilities instead. &
Generally, \textbf{a deposit of 5\% of total liabilities should be treated as large by ALCO}. Where a single depositor breaches the internal single-source concentration limit, the bank must grow its liabilities base to bring the share back within limit, rather than damage the relationship. \\
\addlinespace
\textbf{Undrawn commitments report} &
Aggregate undrawn committed facilities, and the trend over time for both drawn and undrawn committed facilities. &
Aggregate-level; the bank also produces detailed versions. Read as a \emph{contingent} outflow that materialises precisely when the bank can least absorb it. \\
\bottomrule
\end{longtable}
\end{center}

\begin{lrtrapbox}[Two numbers worth locking, and one relationship]
\begin{itemize}
\item \textbf{5\%} of total liabilities is the threshold at which ALCO should treat a single deposit as large. USD 50 million is the illustration, not the rule.
\item The deposit tracker's purpose is the \textbf{LTD (loan-to-deposit) ratio}, not the LCR.
\item The \textbf{funding gap report and the cash flow survival horizon report come from the same underlying report}. An option treating them as independent sources has missed the link.
\end{itemize}
\end{lrtrapbox}

% =====================================================================
\lo{LTR-10 c}{Explain the process of reporting a liquidity stress test and interpret a liquidity stress test report.}

\begin{center}
\small
\begin{tabularx}{\textwidth}{@{}p{4.3cm} X@{}}
\toprule
\rowcolor{FRMSoft}\textbf{Report} & \textbf{What it contains} \\
\midrule
\textbf{1. Cash flow survival report} \newline \emph{(most important)} &
Daily cash inflows and outflows under normal conditions (BAU) \textbf{and} after considering mitigating actions such as selling assets or raising additional funds. It reveals the bank's \textbf{cash flow survival horizon} --- how long the bank can operate. \\
\addlinespace
\textbf{2. Liquidity measures} &
The bank's liquidity situation across categories: funding sources (wholesale, retail and so on), types of assets (marketable, non-marketable), and currency exposure. \\
\addlinespace
\textbf{3. Liability stickiness report} &
How likely depositors are to withdraw their funds during a crisis. \textbf{Stickier liabilities are less likely to be withdrawn.} \\
\addlinespace
\textbf{4. Quarterly line-by-line stress test report} &
The impact of various stress scenarios --- reduced assets or liabilities --- on the bank's liquidity ratios, and \textbf{the probability of each scenario occurring}. \\
\bottomrule
\end{tabularx}
\end{center}
\figcap{The cash flow survival report is named as the most important of the four, and it is the only one that reports the position \emph{twice} --- before and after mitigating actions. That before-and-after structure is what turns a survival horizon into a decision.}

\begin{lrkeybox}[Interpreting the survival horizon]
In the source's worked illustration, mitigating actions extend the survival period from 49 days to \textbf{27 days} in the stressed case --- still below the Basel III requirement, so on the strength of this report \textbf{the bank must take action to address the liquidity shortage}. The interpretation step is the point of the objective: a survival horizon is read against a regulatory minimum, not in isolation.
\end{lrkeybox}

\begin{lrtrapbox}[Stickiness runs the opposite way to intuition for some readers]
\textbf{Sticky} $=$ \emph{less} likely to be withdrawn $=$ \emph{better} for the bank. The word describes the liability's reluctance to leave, not its risk. A statement equating high stickiness with high withdrawal risk is the polarity build, and the word invites it.
\end{lrtrapbox}

\begin{lrtrapbox}[Consolidated trap summary — LTR-10]
\begin{itemize}
\item \textbf{Role.} The \textbf{regulator} stipulates reporting frequency; the \textbf{ALCO} views the reports.
\item \textbf{Scope.} FSA047 and FSA048 move from \textbf{weekly to daily} under stress. The other reports keep a fixed calendar.
\item \textbf{Definition.} 5\% of total liabilities is the ``large depositor'' threshold for ALCO purposes.
\item \textbf{Sibling.} The deposit tracker exists for the \textbf{LTD ratio}; the daily liquidity report exists for the end-of-day position; the mismatch report exists for the \textbf{time-bucket gap} and feeds the survival horizon.
\item \textbf{Polarity.} \textbf{Sticky} liabilities are \emph{less} likely to be withdrawn.
\item \textbf{Scope.} Undrawn commitments are reported \textbf{separately} because they are drawn down precisely under stress.
\item \textbf{Scope.} The cash flow survival report shows the position \textbf{both} before and after mitigating actions, and the horizon is judged against the Basel III requirement.
\item \textbf{Role.} The stress test report's fourth component carries the \textbf{probability of each scenario}, not just the impact.
\end{itemize}
\end{lrtrapbox}
